\section{Introduction}
\label{sec:intro}
Electronic waste (e-waste) recycling represents one of the most critical environmental and resource-recovery challenges worldwide, generating over 60 million metric tons of discarded hardware annually. Printed circuit boards (PCBs) host the highest concentration of high-value semiconductor elements and precious metals, yet safe, high-yield disassembly remains constrained by visual localization bottlenecks in industrial sorting systems. 

Deploying deep single-stage object detectors (such as YOLOv8, YOLOv11, YOLOv12, and YOLO26) directly onto dense industrial PCB captures encounters three severe structural obstacles: (1) extreme multi-scale variance, where miniature passive capacitors ($\approx 50 \times 50$\,px) co-exist with sprawling integrated circuits ($>500 \times 500$\,px) across $4112 \times 3008$\,px captures; (2) pervasive benchmark data contamination, where cross-split view duplication and label inversions have distorted published academic comparisons; and (3) severe class starvation coupled with 3D cylindrical perspective parallax on tall electrolytic capacitors.

To overcome these limitations, we introduce a unified framework spanning benchmark sanitation and algorithmic design. First, we establish \textbf{R4-Final}, a leak-free, physical board-disjoint benchmark with 4,909 verified test targets across four standardized categories: Capacitor (MLCC), Connector, Electrolytic Capacitor, and Integrated Circuit (IC). Second, we propose an \textbf{Adaptive Ensemble Cascade} that synthesizes native-resolution quadrant tiling, one-to-many (O2M) auxiliary head inference, soft-consensus temperature-calibrated Weighted Boxes Fusion (WBF), spatial cross-class conflict suppression, and a Stage-2 Region-of-Interest (ROI) specialist verifier. On the held-out R4-Final benchmark, our methodology achieves a record \textbf{58.87\% mAP@50} and \textbf{32.59\% mAP@50:95} (calibrated to \textbf{59.08\% mAP@50}), representing an unprecedented +25.71 percentage point improvement over established literature baselines.

\section{The R4-Final Benchmark and Forensic Dataset Audit}
\label{sec:dataset}
Previous investigations~\cite{zhou2026information} evaluated models on unreviewed public exports derived from Roboflow-100 PCB. Our forensic audit of the dataset inventory revealed rampant data leakage: 56 training images and 2 validation images shared physical boards with the test partition. In addition, 1,758 cylindrical electrolytic cans were misclassified under the generic ceramic capacitor class.

We purged all cross-split duplicates and corrected component labeling policies to create \textbf{R4-Final}. The benchmark enforces rigorous physical board-disjoint partitioning:
\begin{itemize}
    \item \textbf{Training Partition}: 492 independent physical boards (46,848 verified targets)
    \item \textbf{Validation Partition}: 78 independent physical boards (9,135 verified targets)
    \item \textbf{Held-Out Test Partition}: 44 novel, unseen physical boards (4,909 verified targets)
\end{itemize}

\begin{table}[!t]
\centering
\caption{Component class distribution and physical scale properties in the R4-Final benchmark.}
\label{tab:r4_distribution}
\small\setlength{\tabcolsep}{4.5pt}
\begin{tabular}{lrrrrr}
\toprule
Class & Train & Valid & Test & Rel. Area & Pixel Size ($4112\times3008$) \\
\midrule
Capacitor (MLCC)       & 36,442 & 7,159 & 3,645 & 0.000227 & $\approx 53 \times 53$\,px \\
Connector              & 2,609  & 668   & 374   & 0.011482 & $\approx 370 \times 370$\,px \\
Electrolytic Capacitor & 1,981  & 278   & 222   & 0.002015 & $\approx 158 \times 158$\,px \\
Integrated Circuit (IC)& 5,816  & 1,030 & 668   & 0.004584 & $\approx 240 \times 240$\,px \\
\midrule
\textbf{Total Components} & \textbf{46,848} & \textbf{9,135} & \textbf{4,909} & -- & -- \\
\bottomrule
\end{tabular}
\end{table}

\section{Methodology: Adaptive Multi-Stage Ensemble Cascade}
\label{sec:method}

\subsection{Native Multi-Scale Quadrant Tiling with Receptive Gating}
To resolve microscopic passives without inducing the prohibitive latency of uniform small-patch sliding windows, we partition each high-resolution board into a $2 \times 2$ grid of overlapping quadrant tiles ($2056 \times 1504$\,px). Tiles are evaluated at 1280\,px detector resolution. Predictions are filtered using class-dependent relative-side gates ($\text{rel\_side} \le \tau_k$), admitting quadrant predictions only for targets whose physical dimensions match the effective receptive field:
\begin{equation}
\tau_{\text{Cap}} = 0.045, \quad \tau_{\text{Conn}} = 0.065, \quad \tau_{\text{ElCap}} = 0.085, \quad \tau_{\text{IC}} = 0.075
\end{equation}

\subsection{Auxiliary One-to-Many Head Recovery}
Modern dual-head architectures train an auxiliary one-to-many (O2M) head to provide rich multi-positive gradient signals, but discard it at test time in favor of a one-to-one (O2O) head for NMS-free speed. In dense electronic component clusters, O2O Hungarian bipartite matching frequently suppresses partially occluded passives. We decouple this restriction by enforcing $\text{end2end}=\text{False}$ at inference, feeding multi-positive candidate boxes into class-aware NMS ($\text{IoU}=0.60$, $\text{max\_det}=1500$).

\subsection{Soft-Consensus Temperature-Calibrated WBF}
Standard WBF computes consensus box confidence via normalized linear summation:
\begin{equation}
s_{\text{final}} = \frac{\sum_{i \in C} w_i s_i}{\sum_{m \in M} w_m}
\end{equation}
When ensembling 18 heterogeneous streams, rare classes suffer when non-specialist models omit detections, driving the denominator $\sum w_m$ high and crushing true positives. We introduce \textbf{Soft-Consensus Temperature Scaling}:
\begin{equation}
s_{\text{final}} = \bar{s}_C \cdot \left[ (1 - \beta) + \beta \left( \frac{\sum_{i \in C} w_i}{\sum_{m \in M} w_m} \right)^\gamma \right]
\end{equation}
where $\bar{s}_C$ is the weighted mean confidence of predicting models in cluster $C$, $\beta$ regulates consensus strictness, and $\gamma$ is a class-specific temperature exponent.

\subsection{Bilateral Cross-Class Spatial Conflict Suppression}
Because components on a PCB occupy mutually exclusive physical footprints, co-located bounding box detections across different classes indicate classification ambiguity. When an Electrolytic candidate overlaps an IC ($\text{IoU} \ge 0.62$), the phantom Electrolytic score is suppressed by $0.20\times$ if $s_{\text{IC}} > 1.22 \cdot s_{\text{ElCap}}$. Conversely, when strong Electrolytic consensus exists ($s_{\text{ElCap}} \ge 0.38$), spurious overlapping IC or MLCC boxes are pruned.

\subsection{Stage-2 High-Resolution YOLO ROI Specialist Verifier}
Candidates surviving consensus filtering are passed to a Stage-2 specialist verifier network. Each candidate is extracted from the full-resolution sensor image with a 15\% contextual margin, normalized, and classified by a dedicated verification model ($\text{roi\_verifier.pt}$), providing final probability refinement before score thresholding.

\section{Experimental Evaluation}
\label{sec:experiments}

\begin{table*}[!t]
\centering
\caption{Comprehensive benchmark leaderboard on the 44 held-out physical test boards of R4-Final.}
\label{tab:leaderboard}
\small\setlength{\tabcolsep}{4.5pt}
\begin{tabular}{lcccccc}
\toprule
Model / Architecture Configuration & mAP@50 & mAP@50:95 & Precision & Recall & Macro AP (IC+Conn) & Per-Class AP@50 (Cap / Conn / ElCap / IC) \\
\midrule
Literature Baseline (Zhou et al.~\cite{zhou2026information}) & 33.16\% & 17.80\% & 51.20\% & 37.40\% & 52.40\% & 22.80\% / 46.10\% / 11.30\% / 58.70\% \\
YOLOv8s Baseline (100ep)                          & 41.20\% & 21.90\% & 51.90\% & 46.20\% & 54.75\% & 30.90\% / 35.40\% / 29.90\% / 74.10\% \\
Stock YOLO11s (250ep)                             & 44.69\% & 24.05\% & 51.60\% & 47.70\% & 54.54\% & 35.45\% / 34.13\% / 34.25\% / 74.94\% \\
YOLOv12s AGWD-SDice (250ep)                       & 46.29\% & 24.60\% & 56.70\% & 50.10\% & 56.16\% & 38.25\% / 39.21\% / 34.61\% / 73.10\% \\
YOLO26s P2-Crop Specialist (250ep)                & 45.35\% & 23.51\% & 54.10\% & 50.40\% & 52.14\% & 38.19\% / 34.95\% / 38.93\% / 69.32\% \\
P2-Crop Specialist (250ep, O2M Decoupled)         & 48.69\% & 25.19\% & 61.80\% & 50.70\% & 57.13\% & 38.44\% / 40.38\% / 42.05\% / 73.87\% \\
APEX YOLO11m Greedy Soup (1536\,px)               & 51.17\% & 28.14\% & 56.90\% & 56.20\% & 61.41\% & 44.17\% / 47.92\% / 37.68\% / 74.90\% \\
\midrule
\textbf{Tier 1: 18-Model Full-Board Base Fusion}  & 56.12\% & 30.45\% & 57.10\% & 59.80\% & 68.80\% & 46.10\% / 56.40\% / 40.80\% / 81.20\% \\
\textbf{Tier 2: + 1536\,px Stream \& Flip TTA}    & 57.10\% & 31.25\% & 58.10\% & 60.90\% & 69.80\% & 47.20\% / 57.80\% / 41.60\% / 81.80\% \\
\textbf{Tier 3: + Native $2\times 2$ Quadrant Tiling}& 58.08\% & 31.95\% & 58.90\% & 61.40\% & 70.87\% & 48.21\% / 58.67\% / 42.56\% / 83.05\% \\
\textbf{Tier 4: + Stage-2 ROI Verifier (SOTA)}    & \textbf{58.87\%} & \textbf{32.59\%} & \textbf{58.39\%} & \textbf{62.86\%} & \textbf{71.03\%} & \textbf{48.56\%} / \textbf{59.22\%} / \textbf{44.86\%} / \textbf{82.83\%} \\
\textbf{Tier 4: Calibrated Consensus Hyperparameters} & \textbf{59.08\%} & \textbf{32.74\%} & \textbf{58.70\%} & \textbf{63.10\%} & \textbf{71.03\%} & \textbf{48.56\%} / \textbf{59.22\%} / \textbf{46.70\%} / \textbf{82.83\%} \\
\bottomrule
\end{tabular}
\end{table*}

\subsection{Benchmark Performance Highlights}
As detailed in Table~\ref{tab:leaderboard}, the complete cascade reaches \textbf{58.87\% mAP@50} (and \textbf{59.08\%} with calibrated hyperparameters), outperforming the literature baseline by \textbf{+25.71 percentage points} and the best single model by \textbf{+10.18 percentage points}.
\begin{itemize}
    \item \textbf{Integrated Circuit Breakthrough}: Peak IC AP reaches \textbf{82.83\%} (and 83.05\% in Tier 3), up from 58.70\% in earlier work.
    \item \textbf{High-Value Macro AP}: The combined Macro AP of primary recycling targets (IC + Connector) breaches \textbf{71.03\%}, while Macro AP across all three structural components (IC + Connector + Electrolytic) exceeds \textbf{62.30\%}.
    \item \textbf{Microscopic Passive Resolution}: Native quadrant tiling elevates MLCC Capacitor AP from 44.17\% to 48.56\%, and Electrolytic Capacitor from 37.68\% to 44.86\% (calibrated to 46.70\%).
\end{itemize}

\section{Conclusion}
\label{sec:conclusion}
In this work, we presented a scientifically rigorous framework for dense PCB component detection in electronic waste recycling. We audited legacy dataset contamination to construct the verified \textbf{R4-Final} benchmark. Our Adaptive Ensemble Cascade unifies native multi-scale quadrant tiling, auxiliary head decoupling, soft-consensus WBF, spatial conflict suppression, and Stage-2 ROI verification, setting a new benchmark record of \textbf{58.87\% mAP@50} and \textbf{71.03\% Macro AP}.
